%%
% 第一章：课题背景与实验目标
%%

\chapter{课题背景与实验目标}
\label{chap:background}

\section{课题背景}

代码覆盖率和执行路径分析是软件测试中的基础能力。前者回答“哪些代码被跑到了”，后者回答“程序是怎么走到那里的”。在课堂作业场景里，重点不只是拿到一个覆盖率百分比，而是把追踪、分支识别和结果呈现这几件事自己实现一遍，弄清楚工具背后的工作机制。

Python 生态里已经有成熟的覆盖率工具，但这次作业更强调原理验证和过程可解释性。因此，本文的核心工作不是介绍某个业务项目，而是设计并实现一个面向 Python 程序的覆盖率与路径分析工具，再用真实项目代码来检验它是否能给出有用的测试信息。

\section{作业目标}

根据课程要求，本次作业需要完成以下几项任务：

\begin{itemize}
    \item 记录 Python 程序运行时的执行路径；
    \item 分析顺序、选择、循环等基本控制结构；
    \item 统计代码覆盖情况，并生成可读的执行分支报告；
    \item 提供清晰的代码说明和使用文档，保证工具具备基本的可维护性。
\end{itemize}

围绕这些要求，本文希望回答三个具体问题：第一，能否从静态源代码中提取出需要关注的控制结构；第二，能否在程序运行时稳定记录实际执行路径；第三，能否把两类信息汇总成对测试设计有帮助的报告。

\section{工具方案概述}

本文实现的工具命名为 \texttt{PathTracer}。它采用“静态分析 + 动态追踪”的组合方案：先用 \texttt{ast} 模块扫描目标文件，找出其中的 \texttt{if}、\texttt{for}、\texttt{while}、\texttt{except} 等控制结构；再用 \texttt{sys.settrace} 在程序运行时记录实际执行到的代码行；最后将静态分支信息与动态执行结果对齐，输出覆盖率统计和执行路径报告。

这里需要说明的是，当前版本统计的是“控制结构命中情况”。也就是说，只要某个控制结构所在的起始行在运行时被执行到，就认为该结构被命中。它还不能像成熟工具那样严格区分 \texttt{if} 的真分支和假分支，也不对所有可能路径做穷举。但对本次课程作业来说，这个定义已经足够支撑路径记录、结构识别和报告生成三项核心要求。

\section{PageIndex实验对象说明}

PageIndex 是一个面向长文档处理的索引与问答项目，能够把 PDF、Markdown 等文档整理成层次化结构，并在此基础上完成检索与回答生成。本文并不把 PageIndex 作为研究主体，而是把它当作真实项目样本，用来验证 PathTracer 在实际代码上的覆盖率统计与路径分析能力。换句话说，PageIndex 在本文里的角色是“被测程序”，不是“被介绍的系统”。

从项目结构上看，与本次实验最相关的部分主要有三层：

\begin{itemize}
    \item 索引构建模块：\path{pageindex/page_index.py}，负责解析源文档、抽取层次结构并生成树形索引结果；
    \item 检索与回答模块：\path{pageindex/search.py}，负责列出已建索引文档、筛选相关文档、并发搜索节点、拼接上下文并生成回答；
    \item 命令行入口脚本：\path{agent_pageindex.py}，负责接收用户问题、创建 Agent，并调用 \texttt{search\_pageindex} 工具完成一次完整问答流程。
\end{itemize}

本文将 \path{agent_pageindex.py} 选为主要实验对象，原因有三点。第一，它处在 PageIndex 的入口位置，既能代表真实使用方式，又能把下层检索流程连接进来，因此比单独写一个玩具示例更有说服力。第二，这个脚本的代码规模适中，适合在课程报告里做较细致的路径分析，不至于因为项目过大而让实验描述失去重点。第三，它本身已经包含了多种典型控制结构：\texttt{\_extract\_text} 中有针对 \texttt{None}、字符串、列表等不同输入类型的条件分流；\texttt{\_print\_verbose\_stream} 中有多层 \texttt{for} 循环、消息去重以及工具消息和模型消息的分类处理；\texttt{main} 函数中则包含参数解析、环境变量检查、普通模式与 \texttt{verbose} 模式的路径分叉。

如果进一步看调用链，\path{agent_pageindex.py} 并不是一个孤立的小脚本。它在创建 Agent 后，会把用户问题交给 \texttt{search\_pageindex}；而该工具内部又会继续执行文档目录读取、文档筛选、节点搜索、上下文构造和答案生成等步骤。也就是说，虽然本报告直接统计的目标文件是入口脚本，但这个入口脚本背后连接的是完整的 PageIndex 检索流程，因此实验结果既保留了真实项目的复杂性，又把分析范围控制在课程作业可讨论的规模内。

基于上述考虑，本文对 PageIndex 的介绍只保留与实验设计相关的部分：它是什么、入口脚本做了什么、为什么适合拿来做覆盖率分析。至于文档问答系统本身的业务细节，并不是本次报告要展开的重点。